% !TEX root = tarasov-qasp-zh.tex
\chapter[导体表面附近的电场线如何分布？]{导体表面附近的电场线如何分布？\label{ch-24}}

\begin{dialogue}

\tchr 让我们把某一个导电的物体放入静电场（electrostatic field）中。你们都清楚地知道，处于场中的导体要用一个叫做电容（capacitance）的量来表征。可是你们问过自己没有：为什么我们说导体的电容，却从来不说电介质（dielectric）的电容？

\stda 我从来没有想到过这个问题。

\tchr 你怎样定义孤立导体（isolated conductor）的电容？

\stda 它是为了使导体的电势（potential）升高一个单位而必须传给这个导体的电荷量（quantity of electricity）。
\pnote{这里定义的是孤立导体的电容：电荷量与电势之比（$C=q/\varphi$），“使电势每升高一个单位所需的电荷量”是同一说法。人教版必修第三册第十章第4节《电容器的电容》定义的是电容器的电容 $C=Q/U$（$U$ 为两极板间的电势差）；孤立导体可看作另一极板在无穷远处的特殊情形。孤立导体的电容超出中学要求，物理竞赛与大学普通物理中常见，例如孤立导体球的电容 $C=4\pi\varepsilon_{0}R$。}

\tchr 要注意，你在这里是把电势当作物体的特征来谈论的。可是直到现在，我们一直把电势看作场的特征；作为这样一种特征，它是逐点变化的：电势是场中相应点的坐标的函数。能不能把它说成物体的特征？如果能，那为什么？

\stdb 如果这物体是导体，就可以。事实是：放入静电场中的导体的所有各点都具有相同的电势。导体是等势体（equipotential body）。

\tchr 你的论断以什么为根据？

\stdb 导体中有自由电荷（free charge）。因此，倘若导体任意两点之间存在电势差，这两点之间就会出现电流。这显然是不可能的。

\tchr 完全正确。可以说，当导体被放入静电场时，导体内的自由电荷会重新分布，使得导体内部的电场强度（field intensity）变为零。这实际上就意味着：导体的所有各点（无论是在导体内部还是在它的表面上）都具有同一个电势。导体各点电势处处相同，这使我们能够谈论“导体的电势”。我要指出，电介质中没有自由电荷，因此不可能发生任何电荷的重新分布。顺便问一句：自由电荷在导体中究竟是怎样重新分布的？
\pnote{电介质并非没有电荷，只是没有能自由移动的电荷——其中的电荷被束缚在分子（原子）内部。放入电场后，束缚电荷发生微小移位，介质表面出现极化电荷，这正是电介质影响电容的机制：平行板电容器极板间充满电介质时电容增大，人教版必修第三册第十章第4节用相对介电常数 $\varepsilon_{r}$ 刻画此事。“电介质影响电容”与“电介质自身有电容”是两回事：后者因电介质不是等势体而无从谈起。}

\stdb 它们集中在导体的表面上。导体上任何一个凸出部分的曲率越大，电荷就越密集。电荷密度最大的地方将是尖锐的尖端。
\pnote{电荷只分布在导体外表面、越尖锐处越密集，这是尖端放电与避雷针的原理（人教版必修第三册第九章第4节《静电的防止与利用》）。另外，“曲率越大电荷面密度越大”是定性规则：一般导体的电荷分布由其整体形状及邻近的导体共同决定，并非该处曲率的单一函数。}

\tchr 正是这样。现在清楚了：静电场中的导体是等势体。由此可见，导体的表面是等势面（equipotential surface）。根据这个结论，请告诉我：静电场的电场线（lines of force）在导体表面附近是怎样分布的？

\stdb 既然电场线总是垂直于等势面，它们就必定成直角地“进入”导体表面。
\pnote{原书 lines of force 一词旧译“电力线”，现行人教版教材称“电场线”（必修第三册第九章第3节《电场 电场强度》）；“电场线与等势面垂直、由电势高处指向电势低处”是人教版必修第三册第十章讲等势面时给出的结论。}

\tchr 可惜，考生们常常不知道这一点。画一幅平行板电容器（parallel-plate capacitor）两板之间放有金属球时场的电场线图，这对你应该毫无困难。然而，考生们照例会被这个问题弄得十分困惑。

\stdb 电场线应当以直角趋近电容器的两板和球的表面。这样，电场线的图景就将与\figr{fig-92} 中所示的相似。
\begin{figure}
\centering
\includegraphics[width=0.6\textwidth]{figs/ch-24/fig-092a}
\caption[电场线。]{电场线。\label{fig-92}}
\end{figure}
\tchr 完全正确。我真无法理解，为什么有些考生认为电场线必定要绕过这个球。

现在来考虑下面的问题。\hlred{点电荷（point charge）$+q$ 位于距地面 $r$ 处。它应当在大地中感应出异号的电荷。结果，在电荷与大地之间就产生电的吸引力。求这个力。我建议你们两位都来想一想这道题。}

\stda 大地中感应出的电荷（induced charge）应当与电荷 $+q$ 等量。由此得出，所求的力等于 $q^{2}/r^{2}$。
\pnote{$q^{2}/r^{2}$ 这类写法出自原书静电学部分采用的高斯单位制：那里库仑定律写作 $F=q_{1}q_{2}/r^{2}$（比例系数为 1），点电荷的场强为 $E=q/r^{2}$。人教版必修第三册第九章第2节《库仑定律》采用 SI 记法 $F=kq_{1}q_{2}/r^{2}$，静电力常量 $k=9.0\times10^{9}\ \mathrm{N\cdot m^{2}/C^{2}}$；本题的最终答案在 SI 中应写作 $F=kq^{2}/4r^{2}$。}

\stdb 我不同意这个结论。\smallcaps{学生甲}假定大地中感应出的电荷集中在一点（\figr{fig-93}~\drkgry{(a)} 中的点 $A$）。然而实际上，感应电荷并不集中在一点，而是分布在大地的表面上。因此，我们可以事先断定：所求的力必定小于 $q^{2}/r^{2}$。
\begin{figure}
\centering
\includegraphics[width=0.8\textwidth]{figs/ch-24/fig-093a}
\caption[电场线。]{电场线。\label{fig-93}}
\end{figure}

\tchr 我完全同意你的意见。那么，该怎样着手求电荷与大地之间的吸引力呢？

\stdb 我觉得，我们必须考察电荷与地面之间的场。地面显然是一个等势面。因此，紧靠地面处，场的等势面在形状上必定接近于平面；与此同时，在电荷附近，等势面必定是球面。这样，我们就能画出等势面（更确切地说，这些面的截面）的定性图景。画好之后，就可以按熟知的规则画电场线。\figr{fig-93}~\drkgry{(b)} 中已经这样画了：电场线用实线，等势面的截面用虚线。

\tchr 请继续你的思路。\figr{fig-93}~\drkgry{(b)} 中这幅电场线图没有使你联想到什么吗？

\stdb 当然联想到了。这幅图确实像两个等量异号点电荷电场线的图。我把这一幅画在右边（见\figr{fig-93}~\drkgry{(c)}）。现在一切都十分清楚。两种情形（见\figr{fig-93}~\drkgry{(b)} 与~\drkgry{(c)}）中，电荷 $+q$ 附近场的面貌是相同的。根据式~\eqref{eq-130}，这意味着两种情形中作用在电荷 $+q$ 上的力相同。于是，所求的力等于 $q^{2}/4r^{2}$。
\pnote{这道题的解法即电像法（method of images，又称镜像法）：用在地面下与 $+q$ 对称处的假想电荷 $-q$ 代替大地表面的感应电荷。这样做之所以成立，是因为合成场同样以地面为等势面，满足与原问题相同的边界条件，而静电场的唯一性定理保证满足同一边界条件的场只有一个——“两图中电荷附近场的面貌相同”这一定性判断的严格根据即在于此。镜像电荷与 $+q$ 相距 $2r$，故 $F=q^{2}/(2r)^{2}=q^{2}/4r^{2}$。电像法超出中学范围，属大学电磁学内容，在物理竞赛中常见。}

\tchr 你的推理无懈可击。这道题清楚地表明，场的概念可以异常富有成效。

\end{dialogue}


\cleardoublepage

% NEWTERMS: electrostatic field | 静电场 | 人教版必修第三册第九章；ch-19–23 尚未翻译，译出后首现括注应前移
% NEWTERMS: capacitance (of a conductor) | （导体的）电容 | 此处为孤立导体电容，见评注；人教版必修第三册第十章第4节为电容器的电容
% NEWTERMS: dielectric | 电介质 | 极化与相对介电常数见评注
% NEWTERMS: isolated conductor | 孤立导体 | |
% NEWTERMS: potential | 电势 | 人教版必修第三册第十章；量符号见原文
% NEWTERMS: equipotential body | 等势体 | 静电平衡的导体
% NEWTERMS: free charge | 自由电荷 | |
% NEWTERMS: field intensity | 电场强度 | 即人教版 electric field strength（E）；ch-23 译出后首现括注应前移
% NEWTERMS: charge density | 电荷密度 | 此处指面电荷密度（surface charge density）
% NEWTERMS: equipotential surface | 等势面 | |
% NEWTERMS: lines of force | 电场线 | 旧译“电力线”，人教版称电场线，见评注；ch-23 有该词正式定义，译出后首现括注应前移
% NEWTERMS: parallel-plate capacitor | 平行板电容器 | |
% NEWTERMS: point charge | 点电荷 | 人教版必修第三册第九章第2节；ch-19–23 译出后首现括注应前移
% NEWTERMS: to induce / induced charge | 感应／感应电荷 | 即静电感应（electrostatic induction）
% NEWTERMS: quantity of electricity | 电荷量 | 旧表述
% NEWTERMS: method of images | 电像法 | 仅见于评注层；又称镜像法
